\section{Penerjemahan Pola Bahasa Tingkat Tinggi ke Assembly}
\label{sec:pola-hll}

Bahasa tingkat tinggi menyediakan \code{if}, \code{while}, dan
\code{for}; prosesor 8086 tidak mengenal satupun kata kunci tersebut.
Semua struktur kontrol dibentuk dari tiga bahan dasar: pembanding
yang menyiapkan \eng{flag} (\asm{CMP}/\asm{TEST}), lompatan bersyarat
yang memeriksanya (\asm{Jcc}), dan lompatan tak bersyarat untuk
melompati blok yang tidak dipilih (\asm{JMP}).

\subsection{Pola if}

Bentuk paling sederhana adalah \code{if} tanpa cabang alternatif.

\begin{pseudocode}[caption={Pseudocode pola \code{if}.}, label={lst:pseudo-if}]
jika nilai > 0 maka
    cetak "positif"
akhir jika
\end{pseudocode}

\begin{program}[caption={Penerjemahan pola \code{if} ke assembly.}, label={lst:asm-if}]
    CMP nilai, 0        ; siapkan flag: nilai - 0
    JLE LEWATI          ; bila nilai <= 0, lompati blok
    MOV AH, 09h         ; cetak string "positif"
    LEA DX, pesan_pos
    INT 21h
LEWATI:                 ; lanjutkan program
\end{program}

Kuncinya terletak pada pembalikan syarat: bahasa tingkat tinggi
menuliskan ``kerjakan blok bila syarat benar'', sedangkan assembly
menuliskan ``lompati blok bila syarat tidak benar''.

\subsection{Pola if--else}

Bila terdapat cabang alternatif, diperlukan dua lompatan: satu untuk
memilih cabang, satu untuk melompati cabang lainnya. Alur tersebut
digambarkan pada Gambar~\ref{fig:alur-if-else}.

\begin{figure}[H]
    \centering
    \begin{tikzpicture}[
        font=\small,
        node distance=0.72cm,
        blok/.style={
            draw=obeBlue, line width=0.7pt, fill=obeBlueLight,
            rounded corners=2pt, minimum width=3.6cm,
            minimum height=0.7cm, align=center, inner sep=3pt
        },
        uji/.style={
            draw=obeTeal, line width=0.7pt, fill=obeTealLight,
            diamond, aspect=2.4, minimum width=4.0cm,
            align=center, inner sep=1pt, font=\scriptsize
        },
        panah/.style={-{Stealth[length=2.2mm]}, line width=0.9pt, obeGray}
    ]
        \node[uji] (kondisi) {syarat benar?};
        \node[blok, below left=1.0cm and 1.2cm of kondisi]
            (ya) {blok ``ya''};
        \node[blok, below right=1.0cm and 1.2cm of kondisi]
            (tidak) {blok ``tidak''};
        \node[blok, fill=obeGrayLight, draw=obeGray, below=1.2cm of kondisi]
            (gabung) {lanjutkan program};

        \draw[panah] (kondisi) -- node[left, font=\scriptsize] {ya} (ya);
        \draw[panah] (kondisi) -- node[right, font=\scriptsize] {tidak}
            (tidak);
        \draw[panah] (ya) |- (gabung);
        \draw[panah] (tidak) |- (gabung);
    \end{tikzpicture}
    \caption{Diagram alur percabangan \code{if--else}. Setiap cabang
    diakhiri lompatan ke titik penggabungan.}
    \label{fig:alur-if-else}
\end{figure}

\begin{pseudocode}[caption={Pseudocode pola \code{if--else}.}, label={lst:pseudo-ifelse}]
jika a = b maka
    cetak "sama"
selain itu
    cetak "beda"
akhir jika
\end{pseudocode}

\begin{program}[caption={Penerjemahan pola \code{if--else} ke assembly.}, label={lst:asm-ifelse}]
    MOV AX, a
    CMP AX, b           ; bandingkan a dengan b
    JNE CABANG_BEDA     ; syarat tidak benar -> ke cabang lain
    MOV AH, 09h         ; cabang "sama"
    LEA DX, pesan_sama
    INT 21h
    JMP GABUNG          ; lompati cabang "beda"
CABANG_BEDA:
    MOV AH, 09h         ; cabang "beda"
    LEA DX, pesan_beda
    INT 21h
GABUNG:
\end{program}

\asm{JMP GABUNG} wajib ada; tanpa instruksi itu, setelah mencetak
``sama'' prosesor akan menjatuhi (\eng{fall through}) blok ``beda''
dan mencetak keduanya.

\subsection{Pola while, do--while, dan for}

Ketiga bentuk perulangan hanya berbeda pada letak pengujian dan
penambahan pencacah. Tabel~\ref{tab:pola-hll} merangkum padanannya.

\begin{table}[H]
    \centering
    \caption{Penerjemahan pola struktur kontrol bahasa tingkat tinggi
    ke assembly. Notasi \emph{kond} berarti syarat yang menghasilkan
    \eng{flag} melalui \asm{CMP} atau \asm{TEST}.}
    \label{tab:pola-hll}
    \small
    \begin{tabular}{@{}p{2.6cm}p{5.0cm}p{6.6cm}@{}}
        \toprule
        \textbf{Pola} & \textbf{Bentuk bahasa tingkat tinggi} &
        \textbf{Padanan assembly} \\
        \midrule
        \code{if} &
        \code{if (kond) \{ A \}} &
        \asm{CMP}/\asm{TEST}; lompati \code{A} bila syarat
        \emph{tidak} benar \\
        \addlinespace
        \code{if--else} &
        \code{if (kond) \{ A \} else \{ B \}} &
        uji; lompat ke \code{B} bila salah; \code{A}; \asm{JMP} ke
        gabungan; \code{B} \\
        \addlinespace
        \code{while} &
        \code{while (kond) \{ A \}} &
        \code{uji}: uji di awal; lompat keluar bila salah; \code{A};
        \asm{JMP} kembali ke uji \\
        \addlinespace
        \code{do--while} &
        \code{do \{ A \} while (kond);} &
        \code{badan}: \code{A}; uji di akhir; lompat kembali bila benar
        (dijalankan minimal sekali) \\
        \addlinespace
        \code{for} &
        \code{for (i=0; i<n; i++) \{ A \}} &
        inisialisasi \reg{CX}; \code{ulang}: \code{A}; \asm{INC}/\asm{ADD}
        pencacah; \asm{LOOP} atau \asm{DEC}+\asm{JNZ} \\
        \addlinespace
        \code{switch} &
        \code{switch (x) \{ case 0: \dots \}} &
        indeks dikalikan ukuran entri; \asm{JMP} tak langsung melalui
        tabel alamat \\
        \bottomrule
    \end{tabular}
\end{table}

\code{while} menguji syarat di awal sehingga badan dapat dijalankan
nol kali, sedangkan \code{do--while} menguji di akhir sehingga badan
selalu dijalankan paling sedikit satu kali. Karena itu, \code{while}
memerlukan \asm{JMP} tambahan untuk kembali ke titik pengujian,
sedangkan \code{do--while} cukup memakai \asm{LOOP} atau \asm{JNZ}.

Perulangan \code{for} dengan pencacah turun paling mudah diterjemahkan
memakai \reg{CX} dan \asm{LOOP}, seperti pada Kode~\ref{lst:for-cx}.

\begin{program}[caption={Penerjemahan pola \code{for} dengan pencacah
\reg{CX}.}, label={lst:for-cx}]
    MOV CX, 8           ; for (i = 8; i > 0; i--)
ULANG_FOR:
    ; badan perulangan A
    LOOP ULANG_FOR      ; i--; ulangi selama i <> 0
\end{program}

\begin{infobox}
Bila pencacah harus naik dari nol ke atas, \asm{LOOP} tidak dapat
dipakai langsung karena \asm{LOOP} selalu mengurangi. Gunakan
\asm{INC} pada pencacah terpisah lalu bandingkan dengan batas atas:

\begin{lstlisting}[language=MASM, numbers=none]
    MOV SI, 0           ; i = 0
ULANG_NAIK:
    ; badan perulangan
    INC SI              ; i++
    CMP SI, 8           ; i < 8 ?
    JB  ULANG_NAIK      ; lanjut bila masih di bawah batas
\end{lstlisting}
\end{infobox}

\subsection{Pola switch dengan Tabel Lompatan}

Struktur \code{switch} dengan banyak cabang tidak efisien bila
diterjemahkan sebagai rangkaian \asm{CMP} dan \asm{Jcc} berantai.
Alternatifnya adalah \term{tabel lompatan} (\eng{jump table}): tabel
berisi alamat setiap cabang yang diindeks dengan nomor kasus.

\begin{pseudocode}[caption={Pseudocode pola \code{switch}.}, label={lst:pseudo-switch}]
pilih nomor:
    kasus 0: cetak "nol"
    kasus 1: cetak "satu"
    kasus 2: cetak "dua"
\end{pseudocode}

\begin{program}[caption={Tabel lompatan untuk pola \code{switch}.}, label={lst:asm-switch}]
.DATA
    TABEL DW OFFSET KASUS0, OFFSET KASUS1, OFFSET KASUS2

.CODE
    MOV SI, nomor       ; SI = nomor kasus (0, 1, atau 2)
    CMP SI, 2           ; periksa batas atas
    JA  DI_LUAR_BATAS   ; nomor di luar jangkauan ditolak
    SHL SI, 1           ; setiap entri berukuran 2 byte
    JMP TABEL[SI]       ; lompat ke alamat yang tersimpan pada tabel

KASUS0:
    MOV AH, 09h
    LEA DX, pesan_nol
    INT 21h
    JMP SELESAI

KASUS1:
    MOV AH, 09h
    LEA DX, pesan_satu
    INT 21h
    JMP SELESAI

KASUS2:
    MOV AH, 09h
    LEA DX, pesan_dua
    INT 21h

SELESAI:
\end{program}

Dua hal penting pada Kode~\ref{lst:asm-switch}. Pertama, indeks harus
dikalikan dengan ukuran entri: karena tabel dideklarasikan dengan
\dir{DW}, indeks digeser satu bit ke kiri dengan \asm{SHL} (lihat
Bab~\ref{chap:logika}). Kedua, batas atas \emph{wajib} diperiksa
sebelum pengindeksan; tanpa \lstasm|CMP SI, 2| dan \asm{JA}, nomor
kasus yang salah melompat ke alamat acak. Beberapa assembler menuntut
penulisan \code{JMP WORD PTR TABEL[SI]}.

\subsection{Studi Kasus}

\begin{studikasusbox}
\textbf{Kasus 1: mencari nilai maksimum.} Diberikan larik berisi
delapan bilangan 16-bit. Elemen pertama dianggap kandidat maksimum,
lalu dibandingkan dengan setiap elemen berikutnya dan diperbarui bila
ada yang lebih besar. Perbandingan memakai \asm{CMP} diikuti
\asm{JAE} (tak bertanda). Program lengkapnya disajikan pada
Subbab~\ref{sec:contoh-bab12}; kompleksitasnya $O(n)$.
\end{studikasusbox}

\begin{studikasusbox}
\textbf{Kasus 2: pengurutan gelembung (\eng{bubble sort}).} Larik
diproses berulang kali; pada setiap lintasan, dua elemen bersebelahan
dibandingkan dan ditukar bila urutannya salah, sehingga elemen
terbesar bergerak ke posisi paling kanan. Perulangan bersarang
diterjemahkan menjadi dua pencacah, sehingga nilai \reg{CX} lintasan
luar disimpan sementara di \eng{stack} dengan \asm{PUSH} dan
\asm{POP} (Bab~\ref{chap:prosedur}). Untuk $n$ elemen, algoritme ini
melakukan sekitar $n^2/2$ pembandingan.
\end{studikasusbox}

\begin{studikasusbox}
\textbf{Kasus 3 (pengantar): pencarian biner.} Pada larik yang
\emph{sudah terurut}, batas bawah dan batas atas disimpan pada dua
register penunjuk, dan titik tengah dihitung dengan
\lstasm|MOV AX, bawah|, \lstasm|ADD AX, atas|, lalu
\asm{SHR AX, 1} (Bab~\ref{chap:logika}). Nilai tengah
dibandingkan dengan nilai yang dicari: bila lebih kecil, batas bawah
dipindahkan; bila lebih besar, batas atas dipindahkan. Perulangan
berhenti ketika nilai ditemukan atau kedua batas bertemu. Algoritme
ini memerlukan paling banyak $\log_2 n$ pembandingan, tetapi hanya
sah bila larik terurut \cite{irvine2011}.
\end{studikasusbox}

Seluruh pola pada subbab ini menegaskan bahwa tidak ada struktur
kontrol baru pada tingkat mesin: yang ada hanyalah pembanding,
lompatan bersyarat, dan lompatan tak bersyarat. Kemampuan membaca
rangkaian tersebut merupakan keterampilan inti pada analisis
\eng{disassembly} dan penelusuran program dengan \eng{debugger}
\cite{evans2020}.
